\section{Univariate Polynomial IOP Gadgets}

\subsection{Polynomial identity testing: the algebraic engine}

The starting point is the Schwartz--Zippel--DeMillo--Lipton lemma \cite{DeMilloLipton1978,Schwartz1980,Zippel1979}.

\begin{lemma}[Schwartz-Zippel]
Let $f\in \F[X]$ be a nonzero polynomial of degree at most $d$. If $r\leftarrow \F$ is uniform, then
\[
\Pr[f(r)=0]\le \frac{d}{|\F|}.
\]
The same statement holds for multivariate polynomials when $d$ is replaced by the total degree.
\end{lemma}

This turns polynomial equality into random-point testing:
\[
    f=g \implies f(r)=g(r),
\]
and conversely
\[
    f(r)=g(r) \text{ at a random }r
\]
strongly suggests $f=g$ provided the field is large relative to the degree. In common cryptographic parameter regimes one often has $p\approx 2^{256}$ and $d\le 2^{40}$, so $d/p$ is negligible.

In a compiled proof system, the verifier does not query raw polynomials. Instead, the prover opens committed polynomials at $r$ using the PCS. The IOP verifier checks the algebraic equality on opened values; the PCS verifier checks that these values are consistent with the earlier commitments.

\subsection{Using vanishing polynomials in IOP gadgets}

The algebraic preliminaries introduced the vanishing polynomial
\[
    Z_\Omega(X)=\prod_{a\in\Omega}(X-a).
\]
The gadgets below repeatedly use the same template: to prove that a polynomial $F$ vanishes on a domain $\Omega$, the prover supplies a quotient $q$ satisfying
\[
    F(X)=q(X)Z_\Omega(X).
\]
For multiplicative subgroup domains this is especially efficient because $Z_\Omega(X)=X^k-1$. For smaller subdomains, such as the set of gate-start positions in a simplified Plonk trace, one should use the actual vanishing polynomial of that subdomain; it need not be of the form $X^k-1$.

\subsection{ZeroTest}

\textbf{Goal.} Given a committed polynomial $f$, prove
\[
    f(a)=0\qquad\forall a\in\Omega.
\]

The prover computes
\[
    q(X)=\frac{f(X)}{Z_\Omega(X)}.
\]
The verifier samples $r\leftarrow\F$ and asks for openings of $f(r)$ and $q(r)$, then checks
\[
    f(r)\stackrel{?}{=}q(r)Z_\Omega(r).
\]

Completeness is immediate. For soundness, if $f$ is not divisible by $Z_\Omega$, then
\[
    f(X)-q(X)Z_\Omega(X)
\]
is a nonzero polynomial of controlled degree. By Schwartz-Zippel, it vanishes at a random $r$ with probability at most degree$/|\F|$.

In compiled SNARK form, $q$ is also committed, and the openings at $r$ are verified by the PCS.

\subsection{SumCheck on a subgroup}

\textbf{Goal.} Given a committed polynomial $f$, prove
\[
    \sum_{a\in\Omega}f(a)=0.
\]
For univariate subgroup domains, a convenient way is to construct an accumulator polynomial $s$ satisfying
\[
    s(\omega^0)=f(\omega^0),
\]
\[
    s(\omega^i)=\sum_{j=0}^{i} f(\omega^j)
    \qquad (1\le i\le k-1).
\]
Then
\[
    \sum_{a\in\Omega} f(a)=0
\]
is equivalent to
\[
    s(\omega^{k-1})=0
\]
plus the local recurrence
\[
    s(\omega X)-s(X)-f(\omega X)=0
    \qquad\text{for all }X\in\Omega.
\]
The recurrence is checked by a ZeroTest on the polynomial
\[
    s(\omega X)-s(X)-f(\omega X).
\]
This is the additive analogue of the ProductCheck below.

\subsection{ProductCheck}

ProductCheck is one of the most important univariate gadgets in Plonk-style systems. It converts a global multiplicative statement into a low-degree recurrence.

\subsubsection*{The problem statement}

Let
\[
    \Omega=\set{1,\omega,\omega^2,\ldots,\omega^{k-1}}
\]
be a multiplicative subgroup. Given a committed polynomial $f$, prove
\[
    \prod_{a\in\Omega} f(a)=1.
\]
Directly multiplying all $k$ values would be expensive for the verifier. Instead, the prover commits to a running product polynomial.

\subsubsection*{The running product polynomial}

Define $t\in\polydeg{k}$ by its values on $\Omega$:
\[
    t(1)=f(1),
\]
\[
    t(\omega^s)=\prod_{i=0}^{s}f(\omega^i)
    \qquad s=1,\ldots,k-1.
\]
Thus
\[
    t(\omega)=f(1)f(\omega),
\]
\[
    t(\omega^2)=f(1)f(\omega)f(\omega^2),
\]
and finally
\[
    t(\omega^{k-1})=\prod_{a\in\Omega}f(a).
\]
Therefore the claimed product equals $1$ if
\[
    t(\omega^{k-1})=1.
\]

The local recurrence is
\[
    t(\omega X)=t(X)f(\omega X)
    \qquad\forall X\in\Omega.
\]
The indexing is worth noticing: the recurrence multiplies by $f(\omega X)$, not $f(X)$, because $t(\omega^s)$ includes the value at $\omega^s$.

\subsubsection*{The cyclic last step}

The final point in the subgroup is $X=\omega^{k-1}$. Since $\omega^k=1$, the recurrence at this last point reads
\[
    t(1)=t(\omega^{k-1})f(1).
\]
At first this may look circular, because $t(1)$ is also the first running-product value. But this is exactly why the boundary condition is chosen as $t(\omega^{k-1})=1$. Under that condition, the last recurrence becomes
\[
    f(1)=1\cdot f(1),
\]
so the cyclic wrap-around is consistent. This small indexing point is one of the most common sources of confusion when first learning ProductCheck.

\subsubsection*{Turning the recurrence into a ZeroTest}

Define
\[
    t_1(X)=t(\omega X)-t(X)f(\omega X).
\]
If the recurrence is correct, then
\[
    t_1(X)=0\qquad \forall X\in\Omega.
\]
Hence $Z_\Omega(X)=X^k-1$ divides $t_1(X)$, so the prover can form
\[
    q(X)=\frac{t_1(X)}{X^k-1}.
\]
The verifier samples $r\leftarrow\F$ and asks for openings of
\[
    t(\omega^{k-1}),\quad t(r),\quad t(\omega r),\quad q(r),\quad f(\omega r).
\]
It accepts if
\[
    t(\omega^{k-1})\stackrel{?}{=}1
\]
and
\[
    t(\omega r)-t(r)f(\omega r)
    \stackrel{?}{=}
    q(r)(r^k-1).
\]

This unoptimized ProductCheck protocol uses two new committed polynomials, $t$ and $q$, and five evaluations. The verifier's non-PCS arithmetic is essentially computing $r^k-1$, which costs $O(\log k)$ field operations by repeated squaring.

\subsubsection*{Soundness intuition}

If the prover lies about the product, then at least one of two things must happen:
\begin{enumerate}[leftmargin=2em]
    \item the boundary value $t(\omega^{k-1})$ is not $1$; or
    \item $t$ does not honestly encode the running product, in which case $t_1$ fails to vanish on $\Omega$.
\end{enumerate}
The second case is caught by the ZeroTest. Since the final check is a random polynomial identity test, the cheating probability is bounded by the relevant degree divided by $|\F|$.

\subsection{Rational ProductCheck}

The same idea proves product identities for rational functions. Suppose the goal is
\[
    \prod_{a\in\Omega}\frac{f(a)}{g(a)}=1,
\]
assuming all denominators are nonzero on $\Omega$.

Define
\[
    t(1)=\frac{f(1)}{g(1)},
\]
\[
    t(\omega^s)=\prod_{i=0}^{s}\frac{f(\omega^i)}{g(\omega^i)}
    \qquad s=1,\ldots,k-1.
\]
Then the boundary condition is again
\[
    t(\omega^{k-1})=1,
\]
and the recurrence becomes
\[
    t(\omega X)g(\omega X)=t(X)f(\omega X)
    \qquad\forall X\in\Omega.
\]
Define
\[
    t_1(X)=t(\omega X)g(\omega X)-t(X)f(\omega X).
\]
Then run ZeroTest on $t_1$ over $\Omega$.

This rational version is the bridge from ProductCheck to permutation and wiring arguments.

In randomized permutation arguments, the denominators depend on verifier challenges. A small set of ``bad'' challenges may make a denominator vanish. Protocol descriptions usually exclude those challenges or absorb their probability into the usual Schwartz--Zippel-style soundness error.

\subsection{PermutationCheck via randomized products}

\textbf{Goal.} Given committed polynomials $f$ and $g$, prove that the two lists
\[
    (f(1),f(\omega),\ldots,f(\omega^{k-1}))
\]
and
\[
    (g(1),g(\omega),\ldots,g(\omega^{k-1}))
\]
are the same multiset.

Define the characteristic polynomials of the two multisets:
\[
    \widehat f(Z)=\prod_{a\in\Omega}(Z-f(a)),
\]
\[
    \widehat g(Z)=\prod_{a\in\Omega}(Z-g(a)).
\]
The two lists are permutations of each other iff
\[
    \widehat f(Z)=\widehat g(Z).
\]
By Schwartz-Zippel, it suffices to choose random $r\leftarrow\F$ and prove
\[
    \widehat f(r)=\widehat g(r).
\]
Equivalently,
\[
    \prod_{a\in\Omega}\frac{r-f(a)}{r-g(a)}=1.
\]
This is exactly a rational ProductCheck with
\[
    F(a)=r-f(a),
    \qquad
    G(a)=r-g(a).
\]
This technique is often called Lipton's trick: equality of large multisets is reduced to equality of two random products. Closely related product and shuffle arguments appear in efficient zero-knowledge shuffle proofs \cite{BayerGroth2012}.

\subsection{Prescribed permutation checks}

The previous gadget proves that two lists are equal up to some unknown permutation. Plonk wiring constraints need more: the permutation is prescribed by the circuit.

Let
\[
    W:\Omega\to\Omega
\]
be a known permutation. In protocols it is often encoded by its own low-degree interpolation polynomial on $\Omega$, but that polynomial should not be interpreted as a cheap function to compose inside $g(W(X))$. Given committed polynomials $f$ and $g$, the goal is
\[
    f(y)=g(W(y))\qquad\forall y\in\Omega.
\]
A direct ZeroTest of
\[
    f(X)-g(W(X))
\]
is not efficient, because the composition $g(W(X))$ can have degree about $k^2$ even when $f$, $g$, and $W$ individually have degree $O(k)$. That would destroy the linear/quasilinear prover goal.

The trick is to encode pairs. Consider the two multisets
\[
    \set{(W(a),f(a)):a\in\Omega}
\]
and
\[
    \set{(a,g(a)):a\in\Omega}.
\]
The reason for introducing pairs is conceptual: the first coordinate identifies a wire slot, while the second coordinate records the value stored in that slot. Equality of the pair multisets says that every destination slot receives exactly the value it should receive.

If these multisets are equal, then the value attached to $W(a)$ on the left matches the value attached to the same point on the right, which implies
\[
    f(a)=g(W(a))
\]
for every $a\in\Omega$.

To test equality of pair-multisets, define bivariate polynomials
\[
    \widehat F(X,Y)=\prod_{a\in\Omega}\bigl(X-YW(a)-f(a)\bigr),
\]
\[
    \widehat G(X,Y)=\prod_{a\in\Omega}\bigl(X-Ya-g(a)\bigr).
\]
Because $\F[X,Y]$ is a unique factorization domain, equality of these products of linear factors is equivalent to equality of the corresponding multisets of pairs.

The verifier samples random $r,s\leftarrow\F$ and asks the prover to prove
\[
    \widehat F(r,s)=\widehat G(r,s).
\]
This is equivalent to the rational product identity
\[
    \prod_{a\in\Omega}
    \frac{r-sW(a)-f(a)}{r-sa-g(a)}
    =1.
\]
Again, a rational ProductCheck proves the identity without explicitly constructing the degree-$k$ bivariate products.

